\documentclass[aspectratio=169]{beamer}

\makeatletter
\def\input@path{{shared/}{shared/classes/ugent-beamer/}{shared/styles/timscolours/}{out.nosync/}}
\makeatother

% Load the UGent theme
\usetheme[language=en,faculty=bw,usecolors]{ugent}
\usepackage[ugent]{timscolours}

% \usepackage{svg}
\input{shared/header}

\usepackage{graphicx} % Required for inserting images
\graphicspath{{../shared/classes/ugent-beamer/}{images/}{out.nosync/images/}}
\usepgfplotslibrary{dateplot}
\usepackage{pgfplotstable}
\usepackage{filecontents}
\usepackage{tikz}
\usepackage{tabularx}
\usepackage{enumitem}
\usepackage{diagbox}

% For absolute positioning of logos
\usepackage[absolute,overlay]{textpos}

\titlegraphic{
  \includegraphics[height=1.5cm]{shared/classes/ugent-beamer/fwo_kleur.png}\hspace{1cm}
  % \includegraphics[height=1.5cm]{images/your-logo2.png}
}
\title{Hierarchical Spatio-Temporal State-Space Model for Mosquito Surveillance}
\author{Rita Verstraeten}
\date{\today}

% Custom title page with logos
\addtobeamertemplate{title page}{%
  \begin{textblock*}{8cm}(2.5cm,0.86\paperheight)
    \includegraphics[height=0.75cm]{fwo_kleur.png}
    \hspace{0.4cm}
    \includegraphics[height=0.8cm, trim=9cm 8cm 9cm 8cm, clip]{ITM_logo.jpg}
    \hspace{0.4cm}
    \includegraphics[height=0.9cm]{IPK_logo.jpeg}
  \end{textblock*}
}{}

% Final slide
\newcommand{\finalframe}{
  \begin{frame}
    \vskip-.6cm
    \begin{minipage}[t][.7\textheight][c]{\textwidth}
        \vfill
        \begin{center}
        {\color{white}\fbox{\begin{minipage}{0.45\textwidth}
          \centering
          \textcolor{white}{\large \textbf{Rita Verstraeten}}\\
          \textcolor{white}{\normalsize PhD Student}\\[0.3cm]
          \textcolor{white}{\Large BIONAMIX}\\[0.8cm]
          \textcolor{white}{\footnotesize rita.verstraeten@ugent.be}\\[0.2cm]
          \textcolor{white}{+32 474 40 45 220}\\[0.2cm]
          \textcolor{white}{www.bionamix.ugent.be}
        \end{minipage}}}
        \hspace{2cm}
        \begin{minipage}{0.3\textwidth}
          \centering
          \Huge \textcolor{white}{Thank you!}\\[0.5cm]
          \Large \textcolor{white}{Questions?}
        \end{minipage}
        \end{center}
        \vfill
    \end{minipage}
    \vspace{15pt}
  \end{frame}
}
\setbeamerfont{author}{size=\Large}

\AtBeginSection[]{
  \begin{frame}
    \vfill
    \centering
    \usebeamerfont{title}\insertsectionhead
    \vfill
  \end{frame}
}

\defbeamertemplate*{subsection page}{ugent}[1][]{%
    \vskip-.6cm
    \begin{minipage}[t][.7\textheight][b]{\textwidth}
        \hspace*{.025\textwidth}{\usebeamerfont{title}\usebeamercolor[fg]{title}\insertsectionhead}\\[0.3cm]
        \hspace*{.025\textwidth}{\usebeamerfont{subtitle}\usebeamercolor[fg]{subtitle}\insertsubsectionhead}
    \end{minipage}
    \vspace{15pt}
}
\newcommand{\subsectionframe}{
  \insectionframetrue
  \frame{\subsectionpage}
  \insectionframefalse
}

\begin{document}

% Title slide
\titleframe


% Outline
\begin{frame}{Outline}
  \tableofcontents
\end{frame}

%===============================================
% SECTION 1: INTRODUCTION AND CONTEXT
%===============================================

\section{Introduction and Surveillance Context}

\begin{frame}{Dengue and Mosquito Surveillance}
  \begin{itemize}
    \item \textbf{Dengue transmission:} Transmitted by \textit{Aedes} mosquitoes (particularly \textit{Aedes aegypti})
    \item \textbf{Modelling goals:}
      \begin{itemize}
        \item Estimate true mosquito presence probability in space and time
        \item Account for observation biases in surveillance data
        \item Use entomological model to inform seasonality in mechanistic epi model. $\rightarrow$ \textbf{keep it simple!}
      \end{itemize}
    \item \textbf{Data source:} Entomological surveillance in Cienfuegos, Cuba 2015 - 2019
      \begin{itemize}
        \item Routine household inspections (monthly)
        \item Quality control re-inspections
        \item Reactive inspections triggered by dengue cases
      \end{itemize}
  \end{itemize}
\end{frame}

\begin{frame}{The Data Problem}
  \begin{itemize}
    \item \textbf{Unequal sampling effort:}
      \begin{itemize}
        \item Routine: every household inspected monthly (systematic)
        \item Reactive: additional inspections when dengue cases occur
      \end{itemize}
    \item \textbf{Observation bias:}
      \begin{itemize}
        \item Reactive surveillance inflates observed positivity during outbreaks
        \item Routine and reactive inspections are \textbf{not labeled} in the data
      \end{itemize}
    \item \textbf{Missing data:}
      \begin{itemize}
        \item Absences are not recorded (partial observation)
        \item We only know when mosquitoes \textit{are} detected, not when they \textit{are absent}
      \end{itemize}
  \end{itemize}
\end{frame}

\begin{frame}{Core Problem}
  \begin{center}
    \large
    \textbf{How do we estimate the true mosquito presence probability when surveillance effort depends on dengue cases?}
  \end{center}
  
  \vspace{1.5cm}
  We need to separate:
    \begin{itemize}
    \item \textbf{True ecology:} Actual mosquito dynamics
    \item \textbf{Observation bias:} Changes in detection due to increased effort
    \end{itemize}
\end{frame}

\begin{frame}{Mosquito positivity (Municipaltiy) rate over time}
\centering
  \includegraphics[width=\linewidth]{11_ento_temporal_positivity.png}
\end{frame}

\begin{frame}{Mosquito positivity (Manzanas) rate over time}
\centering
\vspace{0.3cm}
  \includegraphics[width=1.35\textheight]{15_ento_positivity_trend_25_manzanas.png}
\end{frame}

\begin{frame}{Mosquito positivity (CMFs) rate over time}
\centering
\vspace{0.3cm}
  \includegraphics[width=1.35\textheight]{16_ento_positivity_trend_25_cmfs.png}
\end{frame}
%===============================================
% SECTION 2: RESEARCH OBJECTIVE
%===============================================

\section{Research Objective}

\begin{frame}{Objective}
  \textbf{Estimate:}
  \[
    p_{bt} = \Pr(\text{mosquito immature stages present in block } b \text{ at time } t)
  \]
  
  \vspace{1cm}
  
  \textbf{While accounting for:}
  \begin{itemize}
    \item \textbf{Systematic surveillance:} Baseline, covers all households monthly
    \item \textbf{Reactive surveillance:} Additional inspections triggered by dengue cases
    \item \textbf{Environmental drivers:} NDVI, temperature, rainfall, urbanization, etc.
    \item \textbf{Spatial heterogeneity:} Between-block differences
    \item \textbf{Temporal dynamics:} Month-to-month changes and seasonal patterns
  \end{itemize}
\end{frame}

\begin{frame}{Why This Matters}
  \begin{columns}
    \column{0.5\textwidth}
    \textbf{Without correction:}
    \begin{itemize}
      \item Naive analysis $\Rightarrow$ spurious association between cases and mosquitoes
      \item Inflated positivity during outbreaks
      \item Biased estimates of environmental effects
    \end{itemize}
    
    \column{0.5\textwidth}
    \textbf{With correction:}
    \begin{itemize}
      \item Isolate true ecological drivers
    %   \item Evaluate intervention effectiveness
    %   \item Improve surveillance design
    \end{itemize}
  \end{columns}
\end{frame}

%===============================================
% SECTION 3: MODELING APPROACH (BLOCK-LEVEL)
%===============================================

\section{Block-Level Hierarchical Model}

\begin{frame}{Model Overview}
  \textbf{Three-component structure:}
  \begin{itemize}
    \item \textbf{Ecological process:} Latent mosquito presence probability $p_{bt}$
    \item \textbf{Reactive surveillance bias:} Increased detection under reactive surveillance
    \item \textbf{Observation model:} Binomial (or Poisson approximation)
  \end{itemize}
  
  \vspace{1cm}
  
  \textbf{Key insight:} Separate the ecology from the observation process
\end{frame}

\subsection{Observation Data}

\begin{frame}{Available Data Structure}
  For each block $b$ and time $t$:
  
  \begin{center}
    \begin{tabular}{|c|c|}
      \hline
      \textbf{Data} & \textbf{Description} \\
      \hline
      $y_{bt}$ & Mosquito-positive inspection events \\
      $N_b^{HH}$ & Number of households (fixed) \\
      $C_{bt}$ & Observed dengue cases \\
      $X_{bt}$ & Environmental covariates \\
      \hline
    \end{tabular}
  \end{center}
  
  \vspace{1cm}
  
  \textbf{Assumptions:}
  \begin{itemize}
    \item Every household inspected once per month (systematic)
    \item Each dengue case triggers $\kappa$ additional inspections
    \item $\kappa$ is estimated or fixed
  \end{itemize}
\end{frame}

\begin{frame}{Counting Inspection Events}
  \textbf{Total number of inspection events:}
  \begin{equation*}
    n_{bt} = \underbrace{N_b^{HH}}_{\text{systematic}} + \underbrace{\kappa C_{bt}}_{\text{reactive}}
  \end{equation*}
  
  \vspace{1cm}
  
  \begin{itemize}
    \item $N_b^{HH}$ is known and fixed
    \item $C_{bt}$ is observed
    \item $\kappa = 2$ (fixed; each case triggers $\approx 2$ additional inspections)
    \item Assumes: each dengue case originates from a different household
  \end{itemize}
\end{frame}

\subsection{Ecological Process}

\begin{frame}{Latent Ecological Model}
  The \textbf{true} mosquito presence probability:
  \begin{equation*}
    \text{logit}(p_{bt}) = \alpha + X_{bt}\beta + \underbrace{u_b^{GP}}_{\text{spatial}} + \underbrace{v_t^{global} + v_b^{dev}}_{\text{temporal}}
  \end{equation*}

  \begin{itemize}
    \item $\alpha$ = baseline intercept
    \item $X_{bt}\beta$ = environmental effects (distributed lags)
    \item $u_b^{GP}$ = \textbf{spatial GP} with exponential kernel:
      $\text{Cov}(u_b, u_{b'}) = \sigma_{GP}^2\exp(-d_{bb'}/\rho_{GP})$
    \item $v_t^{global}$ = \textbf{city-wide AR(1) trend}: $v_t = \rho v_{t-1} + \sigma_v \varepsilon_t$
    \item $v_b^{dev} \sim \text{Normal}(0, \sigma_{dev})$ = per-block offset from global trend
  \end{itemize}

  \vspace{0.3cm}
  \textbf{Crucial:} Dengue cases do \textbf{NOT} enter this equation!
\end{frame}

\subsection{Distributed Lag Model for Environmental Covariates}

\begin{frame}{Why Lags Matter}
  \textbf{Problem:} Environmental effects are not instantaneous!
  \begin{itemize}
    \item Temperature/rainfall in week 1 may affect larvae development in weeks 2-4
    \item Urbanization (land cover) may have immediate structural effects
    \item Different covariates have different lag structures
  \end{itemize}
  
  \vspace{1cm}
  
  \textbf{Solution:} Distributed Lag Model with \textbf{random walk prior}
\end{frame}

\begin{frame}{Lag Weights with Random Walk Prior}
  \textbf{Lag weight for covariate } $k$ \textbf{ at lag } $\ell$:
  
  \vspace{0.5cm}
  
  \textbf{Prior specification:}
  \begin{align*}
    w_{k,0} &\sim \text{Normal}(0, 0.5) \quad \text{(immediate effect)} \\
    w_{k,\ell} &\sim \text{Normal}(w_{k,\ell-1}, \sigma_w[k]) \quad \ell = 1, \ldots, L \\
    \sigma_w[k] &\sim \text{Exponential}(2)
  \end{align*}
  
  \vspace{1cm}
  
  \textbf{Interpretation:} Each lag weight is centered on the previous lag weight. Encourages \textbf{smooth} trajectories across lags, preventing wild swings.
\end{frame}

\begin{frame}{Lag Structure in Updated Model}
  \textbf{Separate covariates into two types:}
  
  \begin{itemize}
    \item \textbf{Lagged covariates } ($K_L$): Temperature, rainfall, humidity
      \begin{itemize}
        \item Include lags 0 to $L$ (e.g., $L=8$ weeks)
        \item Each follows random walk prior on lag weights
      \end{itemize}
    \item \textbf{Static covariates } ($K_U$): Urban form, land cover
      \begin{itemize}
        \item No temporal variation, simpler prior
        \item $w_k \sim \text{Normal}(0, 0.5)$
      \end{itemize}
  \end{itemize}
\end{frame}

\begin{frame}{Updated Ecological Model with Lags}
  \textbf{Full specification:}
  \begin{equation*}
    \text{logit}(p_{bt}) = \alpha + \underbrace{\sum_{k \in K_L} \sum_{\ell=0}^{L} w_{k,\ell} X_{k,b,t-\ell}}_{\text{lagged}} + \underbrace{\sum_{k \in K_U} w_k X_{k,b}}_{\text{static}} + u_b^{GP} + v_t^{global} + v_b^{dev}
  \end{equation*}

  \begin{itemize}
    \item $X_{k,b,t-\ell}$ = covariate $k$ in block $b$ at lag $\ell$
    \item $w_{k,\ell}$ = estimated lag weight (random walk prior)
    \item $w_k$ = unlagged weight; $\quad$ current covariates: \texttt{total\_rainy\_days}, \texttt{avg\_VPD}, \texttt{precip\_max\_day}, \texttt{mean\_ndvi} (lags 0--1); \texttt{is\_urban}, \texttt{is\_WUI} (unlagged)
  \end{itemize}
\end{frame}

\begin{frame}{Benefits of Random Walk Prior on Lags}
  \begin{itemize}
    \item \textbf{Smoothness:} Realistic lag trajectories (smooth peak, not jagged)
    \item \textbf{Regularization:} Prevents overfitting across many lag coefficients
    \item \textbf{Interpretability:} Easy to visualize when effect peaks
    \item \textbf{Parsimony:} Reduces effective degrees of freedom
    \item \textbf{Identifiability:} Helps distinguish lagged vs.\ unlagged effects
  \end{itemize}
\end{frame}

\subsection{Reactive Surveillance Bias}

\begin{frame}{Modeling Reactive Surveillance Bias}
  When dengue cases occur, detection probability increases:
  \begin{equation*}
    \text{logit}(p_{bt}^R) = \text{logit}(p_{bt}) + \delta_0 + \delta_1 \log(C_{bt})
  \end{equation*}
  
  \vspace{1cm}
  
  \textbf{Parameters:}
  \begin{itemize}
    \item $\delta_0$ = baseline reactive targeting bias
      \begin{itemize}
        \item Positive $\Rightarrow$ easier to find mosquitoes under reactive surveillance
      \end{itemize}
    \item $\delta_1$ = sensitivity to outbreak intensity
      \begin{itemize}
        \item Positive $\Rightarrow$ targeting intensifies with more cases
      \end{itemize}
  \end{itemize}
  
  \vspace{1cm}
  
  \textbf{This affects only observation, not true ecology!}
\end{frame}

\subsection{Mixture of Surveillance Types}

\begin{frame}{Combining Systematic and Reactive Surveillance}
  The \textbf{fraction} of inspections that are reactive:
  \begin{equation*}
    \omega_{bt} = \frac{\kappa C_{bt}}{N_b^{HH} + \kappa C_{bt}}
  \end{equation*}
  
  \vspace{1cm}
  
  The \textbf{effective observation probability} (mixture):
  \begin{equation*}
    \pi_{bt} = \begin{cases}
      p_{bt} & \text{if } C_{bt} = 0 \\
      (1 - \omega_{bt})p_{bt} + \omega_{bt}p_{bt}^R & \text{if } C_{bt} > 0
    \end{cases}
  \end{equation*}
  
  \vspace{1cm}
  
  \textbf{Interpretation:}
  \begin{itemize}
    \item During no-outbreak: only systematic surveillance, probability = $p_{bt}$
    \item During outbreak: weighted mix of unbiased and biased detection
  \end{itemize}
\end{frame}

\subsection{Observation Model}

\begin{frame}{Observation Model: Beta-Binomial}
  \begin{equation*}
    y_{bt} \sim \text{BetaBinomial}(n_{bt},\; \pi_{bt}\cdot\phi,\; (1-\pi_{bt})\cdot\phi)
  \end{equation*}

  \begin{itemize}
    \item Out of $n_{bt}$ inspection events, $y_{bt}$ are positive
    \item $\pi_{bt}$ = mixture of systematic and reactive detection probabilities
    \item $\phi$ = concentration parameter controlling \textbf{overdispersion}
      \begin{itemize}
        \item $\phi \to \infty$: reduces to Binomial
        \item Estimated from zero-case cells: $\phi \approx 23$ (100-block subset)
        \item Empirical variance/binomial-variance ratio $\approx 25$ $\Rightarrow$ overdispersion real
      \end{itemize}
    \item $\kappa = 2$ is fixed $\Rightarrow$ $n_{bt}$ fully determined by data, no approximation needed
  \end{itemize}
\end{frame}

\begin{frame}{Full Model Specification}
  \textbf{Ecological process:}
  \begin{equation*}
    \text{logit}(p_{bt}) = \alpha + X_{bt}\beta + u_b^{GP} + v_t^{global} + v_b^{dev}
  \end{equation*}

  \textbf{Reactive bias:}
  \begin{equation*}
    p_{bt}^R = \text{logit}^{-1}\!\left(\text{logit}(p_{bt}) + \delta_0 + \delta_1 \log C_{bt}\right)
  \end{equation*}

  \textbf{Observation model:}
  \begin{equation*}
    y_{bt} \sim \text{BetaBinomial}(n_{bt},\; \pi_{bt}\phi,\; (1-\pi_{bt})\phi)
  \end{equation*}

  \textbf{Random effects:}
  \begin{itemize}
    \item $u_b^{GP}$: spatial GP, exponential kernel, $\sigma_{GP} \sim \text{Half-Normal}(0,1)$, $\rho_{GP} \sim \text{InvGamma}(3,150)$
    \item $v_t^{global}$: global AR(1), $\sigma_v \sim \text{Exp}(1)$, $\rho \sim \text{Normal}(0.4, 0.2)$
    \item $v_b^{dev}$: per-block offset, $\sigma_{dev} \sim \text{Exp}(2)$
  \end{itemize}
\end{frame}

\section{Strengths and Limitations}

\begin{frame}{Strengths of Block-Level Approach}
  \begin{itemize}
    \item \textbf{Computationally tractable:} Feasible in Stan
    \item \textbf{Direct inference:} Estimates $p_{bt}$ directly
    \item \textbf{Avoids identifiability issues:} Doesn't try to estimate visit probabilities
    \item \textbf{Clear separation:} Ecology vs.\ observation cleanly separated
  \end{itemize}
\end{frame}

\begin{frame}{Limitations}
  \begin{itemize}
    \item \textbf{Mechanistic detail:} Doesn't explain \textit{why} $p_{bt}$ changes
      \begin{itemize}
        \item Colonization vs.\ persistence not separated; control effect not estimable
        \item $\Rightarrow$ addressed by the Occupancy Block Model
      \end{itemize}
    \item \textbf{Assumption on coverage:} Assumes all households visited monthly
    \item \textbf{Clustering:} Assumes each case originates from a different household
    \item \textbf{GP computation:} Exact Cholesky is $\mathcal{O}(B^3)$ per HMC step
      \begin{itemize}
        \item HSGP approximation available ($\mathcal{O}(B \cdot M^2)$) for scaling to all blocks
      \end{itemize}
    \item \textbf{Unlabelled inspections:} Baseline/reactive split relies entirely on $\kappa = 2$
  \end{itemize}
\end{frame}

\begin{frame}{When to Use This Approach}
  \textbf{Best for:}
  \begin{itemize}
    \item Robust estimation of $p_{bt}$ for prediction
    \item Surveillance evaluation and design
    \item Quantifying environmental drivers
    \item When computational efficiency is important
  \end{itemize}
  
  \vspace{1cm}
  
  \textbf{Not ideal for:}
  \begin{itemize}
    \item Understanding mechanistic drivers (colonization vs.\ persistence)
    \item Evaluating specific control interventions
    \item Very detailed mechanistic modeling
  \end{itemize}
\end{frame}

% \begin{frame}{Possible Extensions}
%   \begin{itemize}
%     \item \textbf{Distributed lag models:} Allow lagged effects of covariates
%     \item \textbf{Temporal autocorrelation:} AR(1) or RW1 for $v_t$
%     \item \textbf{Bayesian estimation of $\delta_0, \delta_1$:} Rather than fixing
%     \item \textbf{Negative binomial:} If overdispersion detected
%     \item \textbf{Hurdle models:} Separate zero-inflation from conditional process
%   \end{itemize}
% \end{frame}



%===============================================
% SECTION: GLMM — FREQUENTIST APPROACH
%===============================================

\section{GLMM --- Frequentist Approach}

\begin{frame}{Why a Frequentist GLMM?}
  \textbf{Role in the modelling pipeline:}\vspace{0.3cm}

  \noindent\textbf{1.~Covariate screening:} fast identification of which environmental predictors and lags are associated with mosquito presence\vspace{0.15cm}

  \noindent\textbf{2.~Prior elicitation:} fixed-effect estimates (log-odds scale) inform prior means for the Bayesian model\vspace{0.15cm}

  \noindent\textbf{3.~Variance calibration:} block and temporal variance estimates inform Stan hyperpriors

  \vspace{0.4cm}
  \noindent\textbf{Key difference from the Bayesian model:}\vspace{0.1cm}

  \noindent MLE via \texttt{glmmTMB} (Laplace approximation) --- no priors, fast, no MCMC.
  $\Rightarrow$ run first, feed results into Stan \\
  \vspace{0.4cm}
  \hspace{0.5cm}
  $\Rightarrow$ \textit{is this considered overfitting?}
\end{frame}

\begin{frame}{Data Expansion: Separating Surveillance Types}
  When $C_{bt} > 0$, each original observation is \textbf{split into two rows}:

  \vspace{0.3cm}
  \noindent\textbf{Baseline row} (systematic surveillance):
  \begin{equation*}
    n^{(base)} = \lfloor(1-\omega_{bt})\,n_{bt}\rfloor, \quad
    y^{(base)} = \lfloor y_{bt}(1-\omega_{bt})\rfloor, \quad
    \texttt{reactive\_shift} = 0
  \end{equation*}

  \noindent\textbf{Reactive row} (case-triggered inspections):
  \begin{equation*}
    n^{(react)} = \lfloor\omega_{bt}\,n_{bt}\rfloor, \quad
    y^{(react)} = \lfloor y_{bt}\,\omega_{bt}\rfloor, \quad
    \texttt{reactive\_shift} = \log(1+C_{bt})
  \end{equation*}

  where $\omega_{bt} = \kappa C_{bt} / n_{bt}$

  \vspace{0.3cm}
  \textbf{Caveat:} positives are split proportionally --- assumes equal positivity rate in both surveillance types (discussed in limitations)
\end{frame}

\begin{frame}{GLMM Linear Predictor}
  For each row $i$ (baseline or reactive):
  \begin{equation*}
    \eta_i = \alpha
      + \sum_{k \in \mathcal{L}} \sum_{\ell=0}^{L} \beta_{k\ell}\, x_{k,b(i),t(i)-\ell}
      + \sum_{k \in \mathcal{U}} \beta_k\, x_{k,b(i)}
      + \delta \cdot \texttt{reactive\_shift}_i
      + v_{t(i)} + \phi_{b(i),t(i)}
  \end{equation*}

  \begin{itemize}
    \item $\mathcal{L}$: lagged covariates (lags $\ell = 0,1,2$): \texttt{rainy\_days\_cat}, \texttt{avg\_VPD}, \texttt{precip\_max\_day}
    \item $\mathcal{U}$: unlagged: \texttt{is\_urban}, \texttt{is\_WUI}
    \item $\delta$: log-odds shift for reactive rows
    \item $v_t \sim \mathcal{N}(0,\sigma_v)$: i.i.d.\ month random intercept
    \item $\phi_{b,t}$: AR(1) within-block temporal effect: \texttt{ar1(year\_month + 0 | block)}
  \end{itemize}
\end{frame}

\begin{frame}{GLMM Observation Model and Random Effects}
  \textbf{Likelihood:}
  \begin{equation*}
    y_i \sim \text{Binomial}(n_i,\; \text{logit}^{-1}(\eta_i))
  \end{equation*}
  specified in \texttt{glmmTMB} as \texttt{cbind(y\_i, n\_i - y\_i) \textasciitilde\ \ldots}

  \vspace{0.5cm}
  \textbf{Configurable random effects} (\texttt{cfg} list in \texttt{GLMM.r}):
  \begin{center}
  \small
  \begin{tabular}{lll}
    \hline
    \textbf{Effect} & \textbf{Formula} & \textbf{Current} \\
    \hline
    Block intercept & \texttt{(1 | block)} & off \\
    Month intercept & \texttt{(1 | year\_month)} & on \\
    Within-block AR(1) & \texttt{ar1(time + 0 | block)} & on \\
    Global AR(1) & \texttt{ar1(time + 0 | all)} & off \\
    Spatial exp.\ cov. & \texttt{exp(xy + 0 | spatial)} & off \\
    \hline
  \end{tabular}
  \end{center}
\end{frame}

\begin{frame}{GLMM Outputs and Role}
  \textbf{Outputs:}
  \begin{itemize}
    \item Fitted baseline probability $\hat{p}_{bt}$ (from baseline rows)
    \item Fitted reactive probability $\hat{p}_{bt}^R$ (from reactive rows, when $\kappa > 0$)
    \item Weighted average: $\hat{\pi}_{bt} = (1-\omega_{bt})\hat{p}_{bt} + \omega_{bt}\hat{p}_{bt}^R$
    \item Fixed-effect table with odds ratios
    \item Diagnostic plots: time-series, QQ, residuals, random effects
  \end{itemize}

  \vspace{0.5cm}
  \textbf{Key limitation --- proportional positive split:}
  \begin{itemize}
    \item \textbf{Are we comparing apples to pears by looking at ecological positivity rate ($p_{b,t}$) vs observed positivity rate?}
    \item \textbf{Uncontrolled variance in AR(1) factor}
    \item Assumes reactive and baseline rows have the \emph{same} positivity rate
    \item In reality reactive inspections target high-risk households $\Rightarrow$ positivity likely higher
  \end{itemize}
\end{frame}

\section{Results}
\subsection{GLMM + temporal AR}

\begin{frame}{GLMM Municipality-overall fit}
  \centering
  \vspace{0.2cm}
    \includegraphics[width=\linewidth]{Model_output/probabilities_timeseries_20260324.png}

\end{frame}

\begin{frame}{GLMM per-manzana fit}
  \centering
  \vspace{0.2cm}
    \includegraphics[width=\linewidth]{Model_output/probabilities_timeseries_random_blocks_20260324.png}
\end{frame}

\begin{frame}{GLMM Output coefficients}
  \centering
  \vspace{0.2cm}
    \includegraphics[height=\textheight]{Model_output/Txt_output_20260324.png}
\end{frame}


\begin{frame}{GLMM QQ plot}
  \centering
  \vspace{0.2cm}
    \includegraphics[height=\textheight]{Model_output/glmm_qqplot_observed_vs_expected_20260324.png}
\end{frame}


\subsection{GLMM + temporal RE + spatial RE}

\begin{frame}{GLMM Municipality-overall fit}
  \centering
  \vspace{0.2cm}
    \includegraphics[width=\linewidth]{Model_output/probabilities_timeseries_20260313.png}

\end{frame}

\begin{frame}{GLMM per-manzana fit}
  \centering
  \vspace{0.2cm}
    \includegraphics[width=\linewidth]{Model_output/probabilities_timeseries_random_blocks_20260313.png}
\end{frame}

\begin{frame}{GLMM Output coefficients}
  \centering
  \vspace{0.2cm}
    \includegraphics[height=\textheight]{Model_output/Txt_output_20260313.png}
\end{frame}


\begin{frame}{GLMM QQ plot}
  \centering
  \vspace{0.2cm}
    \includegraphics[height=\textheight]{Model_output/glmm_qqplot_observed_vs_expected_20260313.png}
\end{frame}

\subsection{GLMM + spatial AC}

\begin{frame}{GLMM Municipality-overall fit}
  \centering
  \vspace{0.2cm}
    \includegraphics[width=\linewidth]{Model_output/probabilities_timeseries_20260314.png}

\end{frame}

\begin{frame}{GLMM per-manzana fit}
  \centering
  \vspace{0.2cm}
    \includegraphics[width=\linewidth]{Model_output/probabilities_timeseries_random_blocks_20260314.png}
\end{frame}

\begin{frame}{GLMM Output coefficients}
  \begin{columns}
    \begin{column}{0.5\textwidth}
      \vspace{0.2cm}
      \includegraphics[height=0.95\linewidth]{Model_output/Txt_output_20260314_part1.png}
    \end{column}
    \begin{column}{0.5\textwidth}
      \vspace{0.2cm}
      \includegraphics[height=0.95\linewidth]{Model_output/Txt_output_20260314_part2.png}
    \end{column}
  \end{columns}
\end{frame}

\begin{frame}{GLMM QQ plot}
  \centering
  \vspace{0.2cm}
    \includegraphics[height=\textheight]{Model_output/glmm_qqplot_observed_vs_expected_20260314.png}
\end{frame}

\section{Questions}
\begin{frame}{Questions for Stijn}
\begin{itemize}
  \item Given that my \textbf{outcome is a rare event} (with low levels when present) (±1\% prevalence) with a reactive surveillance process, is a GLMM with a split-row design (baseline vs reactive rows) the right frequentist approach, or would another model be better? 
  \item Given the data, should I rather tend towards \textbf{cleaning up the data} or should I adjust my modelling strategy? 
  \item What can I actually use the GLMM model for given the extremely high variance on the temporal autoregressive factor? 
\end{itemize}
\end{frame}


\finalframe

\end{document}
